\subsection{疲劳分析}
\label{sec:fatigue-4in1}

\subsubsection{取值}

四合一双功能刮管器承力件材料为 4330V-155KSI（$\sigma_s\ge930$\,MPa、$\sigma_u\ge1\,033$\,MPa）。按第~\ref{sec:fatigue-3in1}~节的方法，取 $k_a=4.51\times1\,033^{-0.265}=0.717$（$Ra1.6$ 机加工面）、$k_b=0.748$（危险截面 $\phi88$\,mm）、$k_c=0.85$（轴向拉压）、$k_d=1.0$、$k_e=0.814$（99\% 可靠度），五项系数之积为 0.371，得 $1\times10^{6}$ 次处的疲劳极限 $S_e=0.5\times1\,033\times0.371\approx191.5$\,MPa。$1\times10^{3}$ 次处取 $0.9\sigma_u=929.7$\,MPa，该值略低于屈服强度 930\,MPa，无需截断。由五点拟合得 $\sigma_a=4\,512.7\,N^{-0.2287}$，$1\times10^{3}$ 次处的许用应力幅 $\sigma_e=929.7$\,MPa。

评定应力同样按剔除奇异点、取梯度平缓区的方式读取。轴向拉伸工况的峰值 923.1\,MPa 位于下心轴与心轴连接部位的变截面台阶根部，属几何尖角引起的数值奇异；距该点约 3\,mm 处的两个节点读数为 588.26\,MPa 与 584.39\,MPa，二者相差 0.7\%，取大值 588.3\,MPa 作为评定应力。抗扭工况的峰值 965.1\,MPa 位于孔口过渡尖角，探针读数为 570.99\,MPa 与 572.49\,MPa，取 572.5\,MPa。两个工况均按 $R=0$ 脉动循环评定，因起下钻提拉与上卸扣扭矩均自零加至峰值，最低载荷为零，故 $\sigma_a=\sigma_m=\sigma_{\max}/2$。评定结果列于表~\ref{tab:fatigue_4in1}。

\begin{table}[htbp]
\centering
\caption{四合一双功能刮管器的疲劳评定结果（$R=0$，200 趟）}
\label{tab:fatigue_4in1}
\begin{tabularx}{\textwidth}{lCCC}
\hline
项目 & 轴向拉伸 & 抗扭 & 判据 \\
\hline
额定工况 & 1\,000\,kN & 17\,kN$\cdot$m & --- \\
评定应力 / MPa & 588.3 & 572.5 & --- \\
应力幅与平均应力 / MPa & 294.1 & 286.3 & --- \\
等效对称应力幅 / MPa & 411.2 & 395.5 & --- \\
计算寿命 / 次 & $3.54\times10^{4}$ & $3.86\times10^{4}$ & --- \\
200 次累积损伤 & $5.96\times10^{-3}$ & $5.18\times10^{-3}$ & $\le 1.0$ \\
疲劳安全系数 & 1.66 & 1.71 & $\ge 1.5$ \\
静强度安全系数 & 1.58 & 1.62 & $\ge 1.5$ \\
寿命裕度 & 约 177 & 约 193 & --- \\
\hline
\end{tabularx}
\end{table}

将评定应力代入计算流程即可复核表中数值。拉伸工况：$\sigma_{a,\mathrm{eq}}=294.1/(1-294.1/1\,033)=411.2$\,MPa，反解得 $N=(4\,512.7/411.2)^{1/0.2287}=3.54\times10^{4}$ 次，与寿命云图在该截面的探针读数 35\,408 次一致；$n_f=1/(294.1/929.7+294.1/1\,033)=1.66$，$n_s=930/588.3=1.58$。抗扭工况的评定应力与拉伸接近，亦按脉动循环取值；寿命 $3.86\times10^{4}$ 次取自云图在该节点的读数 38\,600 次，较按 395.5\,MPa 反算的 $4.2\times10^{4}$ 次低 8\%，属第~\ref{ss:fatigue-probe}~节所述的同源差异，取云图读数作为保守估计。

累积损伤按云图损伤读数折算，而非由探针寿命直接相除。拉伸工况该截面的单次循环损伤为 $2.98\times10^{-5}$（对应寿命 $3.36\times10^{4}$ 次，较探针读数低 5\%），200 趟累积损伤为 $5.96\times10^{-3}$；抗扭工况单次损伤为 $2.59\times10^{-5}$，200 趟为 $5.18\times10^{-3}$。两种算法相差 5\% 以内，取寿命较低的一侧，结论偏于安全。

两个工况的静强度安全系数（1.58、1.62）与疲劳安全系数（1.66、1.71）均高于 1.5，累积损伤较判据 1.0 低两个量级，寿命裕度分别约为 177 倍与 193 倍，均满足要求。抗扭工况并不比拉伸工况更严酷：两者的评定应力仅相差 2.7\%，应力幅相同，寿命自然接近；仅当扭矩按完全反向（$R=-1$）取定时，应力幅才会翻倍。按 $R=-1$ 折算，抗扭寿命由 $3.86\times10^{4}$ 次降至约 $8.3\times10^{3}$ 次，200 趟累积损伤升至约 $2.4\times10^{-2}$，仍满足判据，但裕度仅剩四十余倍，在与其他工具横向比较时应予注意。

按 API RP 7G 安全系数 1.6 反推的 625\,kN 工作载荷复核，评定应力降至 367.7\,MPa，等效对称应力幅为 223.7\,MPa，寿命提高至 $5.06\times10^{5}$ 次，200 趟累积损伤为 $3.95\times10^{-4}$，静强度与疲劳安全系数分别提高至 2.53 与 2.66。服役载荷下的裕度较额定工况宽一个量级，表明 1.58 的静强度安全系数贴近判据系额定载荷的定义所致，而非材料或结构的强度不足。

控制部位位于螺纹接头附近的变截面过渡区，与本章强度校核的结论一致。加工时应保证下心轴与心轴连接处的过渡圆角不小于 R3，避免尖角成为疲劳裂纹源；成品螺纹根部应按 NB/T\,47013.4 进行磁粉探伤，以发现加工刀痕与磨削烧伤。
